\scanpage{252}
\mspage{126}
% archive sheet number: 251

Notons ceci :

% slanted
\margin{Cf version plus forte dans 9.\unsure{8|9}, plus bas (débarrassée de l'hyp. Loc (6)).}

\uline{Proposition 9.7}. Supposons \deleted{l'axiome} que $W$ satisfait $\mathrm{Loc}\,(1,2,3,\deleted{\unsure{6}}5)$, i.e. c'est un localiseur fondamental satisfaisant la condition de décomposition (forme faible). \deleted{\unsure{Le}} Si l'axiome Loc $(8)$ de connexité n'est \uline{pas} satisfait, alors \deleted{tous les \unsure{éléments} \illegible{1} \unsure{flèches}}\inserted{la sous-catégor. $\mathrm{Hot}^{\bullet}_W$} des $X\in\mathrm{Cat}$ non vides \unsure{est|\deleted{est}} équiv. à la cat. ponctuelle (donc $\mathrm{Hot}_W$ est équivalente soit, à la catégorie ponctuelle, \deleted{soit à cat} si

% slanted
\margin{\deleted{\unsure{Il suffit de savoir que la condition $2^\circ$ de b) soit satisfaite}}}

% slanted
\inserted{$\varnothing\to e$ est un iso dans $\mathrm{Hot}_W$, soit à $\Delta^1$ dans le cas contraire}

% slanted
\margin{Le premier cas est aussi celui où $\exists f\in W$ avec $\pi_0(f)$ \uline{non surjectif}, ou encore $\varnothing\to X$ dans $W$, avec $X$ non vide.}

Ceci \deleted{\unsure{généralise}}\inserted{est une variante du} th.~2 de \uline{XII} (p.~20), \deleted{\unsure{auquel on}} où, en faisant l'hypothèse Loc $(5c')$, plus forte que Loc $(\unsure{7|5})$ (et $=$\inserted{pour} Loc $(7\,\mathrm{bis})$), mais avec la conclusion plus forte que $W=\mathrm{Fl}(\mathrm{Cat})$

\uline{Dém} Soit
\[
 f:X\unsure{-|\to}Y,\quad f\in W,
\]
telle que
\[
 \pi_0(f)\quad\text{ne soit pas \uline{surjective}.}
\]

\deleted{Alors par Loc} Si $Y_0$ est une comp. connexe qui n'est pas dans $\mathrm{Im}\,\pi_0(f)$, alors $X_0=X\mid Y_0$ est vide, et \deleted{\unsure{par}}\unsure{par} Loc $(7)$ $\varnothing\to Y_0$ est dans $W$. Donc $\forall Z$ de $\mathrm{Cat}$, $\underset{\substack{\|\\\varnothing}}{\varnothing\times Z}\longrightarrow Y_0\times Z$ est dans $W$, \deleted{Pour toute}\inserted{et pour} \deleted{\unsure{toute}} flèche $g:Z-Z'$, $\mathrm{id}_{Y_0}\times g:Y_0\times Z\to Y_0\times Z'$ est dans $W$. Prenant $\unsure{g'|g}$ telle que $\pi_0(g)$ ne soit pas injective, p.ex. $e\amalg e\to e$, on \unsure{trouve|obtient} une flèche $\deleted{\unsure{g}}\unsure{g'|f}=\mathrm{id}_{Y_0}\times g$ \inserted{dans $W$}, telle que $\pi_0(\unsure{g|f})$ ne soit pas injective.
